\section{Neural Networks with Physical Structure}



\begin{frame}{Informed UDE}
    
        
    \begin{block}{\textbf{Exploiting All Physical Knowledge}:}
        \begin{itemize}
            \item Mass matrix $\mathbf{M}$: \textbf{Positive definite} and \textbf{symmetric}.
            \item Friction: \textbf{Positive definite}, \textbf{diagonal/constant} matrix.
            \item Last term in $\mathbf{M}(\mathbf{q})$ matrix is constant as well as $\mathbf{d_{11}}$ and $\mathbf{d_{22}}$ in $\mathbf{D}$ which is diagonal.
            \item \textbf{Implementation}: Use of \textbf{Cholesky decomposition} to guarantee $\mathbf{M}$ positive definite.
        \end{itemize}
    \end{block}
    \vspace{1em}
    \begin{figure}
        \centering
        \includegraphics[width=0.8\textwidth]{08_informed_predictors/Dyn_Informed_UDE.pdf}
    \end{figure}
    
\end{frame}

\begin{frame} {Informed UDE - Results}
    \begin{itemize}      
        \item \textbf{Implementation}: Use of \textbf{Cholesky decomposition} to guarantee $\mathbf{M}$ positive definite.
        \item \textbf{Better} takes less to converge, more stable.
        \item \textbf{Troubles} in estimating entries in highly chaotic regions.
    \end{itemize}

    \textbf{Can we plug in even more knowledge?}
\end{frame}

\begin{frame} {"More" Informed UDE}
    \begin{block}{Idea}      
        If we somehow could give the matrix M to the system, so as to avoid inversion problems...
    \end{block}

    \begin{figure}
        \centering
        \includegraphics[width=0.8\textwidth]{08_informed_predictors/Dyn_More_informed.pdf}
    \end{figure}
\end{frame}

\begin{frame} {Informed UDE - Drawing Conclusions}
    \begin{itemize}      
        \item \textbf{Better} takes less to converge, more stable.
        \item \textbf{Troubles} in estimating entries in highly chaotic regions.
    \end{itemize}
    \vspace{1em}
    \begin{table}[h]
        \centering
        \caption{Comparison of the three UDE methods}
        \label{tab: metodi}
        \begin{tabular}{ l c c p{2.5cm} }
            \toprule
            \textbf{Metodo} & \textbf{Time to converge} & \textbf{Informations} & \textbf{Comments} \\
            \midrule
                Uninformed      & High                  & None                       & Unstable, Dumb    \\
                Informed      & Medium                 & Structural properties           & Stable, Complicated Identification             \\
                More Informed       & Medium                 & M matrix                    & Easier Identification, Unrealistic       \\
            \bottomrule
        \end{tabular}
    \end{table}
\end{frame}

% \begin{frame}
%     \frametitle{Informed UDE for Matrices Entries Estimation}

%     We can enhance the UDE approach `informing' better our UDE

%     \begin{figure}
%         \centering
%         \includegraphics[width=0.4\textwidth]{images/informed_UDE.png}
%     \end{figure}

%      \begin{figure}
%         \centering
%         \includegraphics[width=0.4\textwidth]{images/better_informed_UDE.png}
%     \end{figure}

% \end{frame}


% \begin{frame}
%     \frametitle{Results for Physical Parameters Estimation}

%     Learning directly on the entries makes the task easy.

%     Gauge invariance issue can affect the estimation using UDE. 

%     Informing better UDE helps a lot. 

%     \begin{figure}
%         \centering
%         \includegraphics[width=0.7\textwidth]{images/matirces_entries_results.png}
%     \end{figure}
% \end{frame}





\begin{frame} {Structured Lagrangian Neural Network (SLNN) - Architecture (1)}
    
        Instead of learning a black box $\mathcal{L}$ it is possible to expolit known physical properties of mechanical systems. 

        $$ \mathcal{L} = T - V = \frac{1}{2} \dot{q}^{\top} M \dot {q} - V$$
        \begin{itemize}
            \item  To learn a positive definite $M$ it is possible to use \textbf{Cholesky decomposition}: 
                    $$M = L L^{\top}$$

                    Where $L$ is a learned \textbf{lower triangular} matrix and \textbf{softplus} is applied to its diagonal entries. 
            
            \item The potential energy $V$ is learned via an MPL that has as inputs the state
            \item Additional: \textbf{friction} matrix learned and rendered positive definite via softplus

        \end{itemize}
               
\end{frame}

\begin{frame} {Structured Lagrangian Neural Network (SLNN) - Architecture (2)}
    \begin{figure}[h]
        \centering
        \includegraphics[width=0.5\linewidth]{stock.png} \\
        \vspace{-1em}
    \end{figure}

\end{frame}


\begin{frame}{Structured Lagrangian Neural Network (SLNN) - Comparison LNN vs SLNN vs Baseline}
    
    % Define commands (It is better to put these in the preamble, but they work here)
    \newcommand{\cmark}{\textcolor{green!70!black}{\ding{51}}} 
    \newcommand{\xmark}{\textcolor{red}{\ding{55}}}             

    % This makes the \thead headers bold automatically
    \renewcommand\theadfont{\bfseries}

    \centering
        % 4 columns: 1 for Metric name, 3 for the models
        \begin{tabular}{cccc}
            \toprule
            \textbf{Metric} & \textbf{Baseline} & \textbf{LNN} & \textbf{SLNN} \\
            \midrule
            Energy conservation         & \xmark & \cmark & \cmark \\
            Friction                    & \cmark & \xmark & \cmark \\
            Broken physics sensitivity  & \xmark & \cmark & \cmark \\
            Stable                      & \cmark & \xmark & \cmark \\
            Not-only mechanical systems & \cmark & \cmark & \xmark \\
            \bottomrule
    \end{tabular}

    \vspace{3em}

    It's possible to add different \textbf{modules} to SLNN such as friction to learn different models. 

\end{frame}


